\documentclass[11pt]{article}

\usepackage{amsmath,amssymb,amsthm,mathtools}
\usepackage[T1]{fontenc}
\usepackage{lmodern}
\usepackage[margin=1in]{geometry}
\usepackage{enumitem}
\usepackage{booktabs}
\usepackage{hyperref}
\usepackage{microtype}
\usepackage{xcolor}

\hypersetup{
  colorlinks=true,
  linkcolor=blue,
  citecolor=blue,
  urlcolor=blue
}

\newtheorem{theorem}{Theorem}[section]
\newtheorem{proposition}[theorem]{Proposition}
\newtheorem{lemma}[theorem]{Lemma}
\newtheorem{corollary}[theorem]{Corollary}
\newtheorem{remark}[theorem]{Remark}
\newtheorem{definition}[theorem]{Definition}
\newtheorem{placeholder}[theorem]{Placeholder}

\newcommand{\Z}{\mathbb Z}
\newcommand{\HLone}{H_{L1}}

\title{D5 129: the front-end channel reduction package for $M1$}
\author{}
\date{March 2026}

\begin{document}
\maketitle

\begin{abstract}
This note revisits the remaining theorem-side placeholder $M1$ from the
revised consolidated manuscript.
The main point is that the old wording ``channel reduction from the full
odd-$m$, $d=5$ problem to the best-seed channel'' mixed two logically
separate front-end tasks:
\begin{enumerate}[label=\textup{(\roman*)},leftmargin=2.5em]
\item isolate the best endpoint seed; and
\item for that seed, prove that the defect-template quotient leaves only one
      unresolved channel.
\end{enumerate}
The second implication is already theorem-level once one accepts the promoted
best-seed defect package of the existing RoundY front end.
So inside the current accepted front-end package, $M1$ is closed.
If one wants a fully first-principles manuscript-order proof, the remaining
externality is now much smaller than the old placeholder: it is exactly the
promotion of the front-end seed-isolation and defect-quotient theorems usually
referred to as \textup{032} and \textup{033}.
\end{abstract}

\section{The old $M1$ placeholder and its real content}

The revised consolidated manuscript leaves the following theorem-side
placeholder open:
\begin{quote}
prove a channel reduction from the full odd-$m$, $d=5$ problem to the
best-seed channel
\[
R1 \to \HLone.
\]
\end{quote}

As stated, this hides two different front-end reductions.
First one must know which endpoint seed is the correct live seed.
Only after that does it make sense to ask which defect channels remain inside
that seed.

The existing RoundY record already separates those roles informally:
\begin{enumerate}[label=\textup{(F\arabic*)},leftmargin=2.8em]
\item the front-end seed isolation usually referred to as \textup{032}:
      the best endpoint seed is
      \[
      L=[2,2,1], \qquad R=[1,4,4];
      \]
\item the best-seed defect quotient usually referred to as \textup{033}:
      for that seed, the balanced Latin defect quotients to four overfull
      families against four hole families, three channels repair directly, and
      only one unresolved channel remains.
\end{enumerate}

The purpose of this note is to package this split honestly and to record the
exact theorem that $M1$ already gets from the current front-end package.

\section{Accepted front-end inputs}

We write the best endpoint seed as
\[
L=[2,2,1], \qquad R=[1,4,4].
\]

\begin{definition}[front-end seed package]
\label{def:F032}
The \emph{front-end seed package} is the statement that in the odd-$m$, $d=5$
endpoint-local problem, after the earlier tiny static repair branches are
pruned, it is enough to work with the best endpoint seed
\[
L=[2,2,1], \qquad R=[1,4,4].
\]
We refer to this as the \emph{032 package}.
\end{definition}

\begin{definition}[best-seed defect package]
\label{def:F033}
The \emph{best-seed defect package} is the following theorem bundle for the
seed $[2,2,1]/[1,4,4]$ at odd modulus $m$.
\begin{enumerate}[label=\textup{(\alph*)},leftmargin=2.5em]
\item Per color, the balanced Latin defect quotients to four overfull families
      and four hole families.
\item The changed source families are exactly
      \[
      L3,\ R1,\ R2,\ R3.
      \]
\item The hole families are exactly
      \[
      H_{L3},\ \HLone,\ H_{R2},\ H_{R3}.
      \]
\item The channels $L3$, $R2$, and $R3$ repair directly.
\item The only unresolved channel is
      \[
      R1 \to \HLone.
      \]
\item The target hole family of that unresolved channel is exactly
      \[
      \HLone
      =\{(q,w,u,\Theta)=(m-1,m-1,u,2): u\neq 2\}.
      \]
\end{enumerate}
We refer to this as the \emph{033 package}.
\end{definition}

\begin{remark}
Items \textup{(a)}--\textup{(f)} are exactly the front-end statements already
frozen informally by the existing RoundY artifact and README layer:
\begin{enumerate}[label=\textup{(\roman*)},leftmargin=2.5em]
\item the best endpoint seed is the seed $[2,2,1]/[1,4,4]$;
\item its balanced Latin defect is $10m^2$ on the pilot range;
\item per color the defect quotient has four overfull and four hole families;
\item the direct repairs are $L3$, $R2$, and $R3$;
\item the unique unresolved channel is $R1\to \HLone$;
\item the hole family $\HLone$ has the exact formula above.
\end{enumerate}
The present note does not attempt a first-principles reconstruction of the
front-end compute-backed proofs behind those statements; it packages what $M1$
gets once those front-end statements are accepted.
\end{remark}

\section{The exact theorem $M1$ needs}

The theorem-side graph manuscript does not need the whole front-end search
history.
It only needs the following reduction statement.

\begin{theorem}[channel reduction from the accepted front-end package]
\label{thm:M1-from-front-end}
Assume the front-end seed package \textup{032}
(Definition~\ref{def:F032}) and the best-seed defect package \textup{033}
(Definition~\ref{def:F033}).
Then every odd-$m$, $d=5$ instance is either already closed by an earlier
front-end repair branch or reduces to the best-seed channel
\[
R1 \to \HLone.
\]
Moreover the target hole family of that channel is exactly
\[
\HLone
=\{(q,w,u,\Theta)=(m-1,m-1,u,2): u\neq 2\}.
\]
\end{theorem}

\begin{proof}
By the front-end seed package \textup{032}, the full odd-$m$, $d=5$
endpoint-local front end reduces to the single seed $[2,2,1]/[1,4,4]$.
Inside that seed, the best-seed defect package \textup{033} says that the
changed source families are exactly
\[
L3,\ R1,\ R2,\ R3,
\]
that the channels $L3$, $R2$, and $R3$ repair directly, and that the only
unresolved channel is
\[
R1 \to \HLone.
\]
So after the earlier front-end repair branches and the direct repairs are
removed, the unique remaining branch is the best-seed channel
$R1\to \HLone$.
The exact formula for the target hole family is part of the same
\textup{033} package, so the reduction carries that formula as well.
\end{proof}

\begin{corollary}[the theorem-side content of $M1$]
\label{cor:M1-closed-in-accepted-front-end-package}
Inside the current accepted front-end package, $M1$ is closed.
Equivalently, the graph-theoretic route may now begin with the best-seed
channel
\[
R1 \to \HLone
\]
as a theorem-level starting point, provided the manuscript states explicitly
that it is invoking the accepted front-end packages \textup{032} and
\textup{033}.
\end{corollary}

\begin{proof}
This is exactly Theorem~\ref{thm:M1-from-front-end}.
\end{proof}

\section{What is still external if one wants a fully independent manuscript}

The previous theorem closes $M1$ only after one accepts the front-end packages
\textup{032} and \textup{033} as theorem-level inputs.
So the old single $M1$ placeholder can now be replaced by a much sharper pair
of front-end placeholders.

\begin{placeholder}[front-end seed isolation]
\label{ph:032}
Promote the best-endpoint-seed statement usually referred to as
\textup{032} into a manuscript-order theorem with explicit odd-$m$ quantifier:
prove that, after the earlier tiny static repair branches are pruned, the
odd-$m$, $d=5$ endpoint-local front end reduces to the seed
$[2,2,1]/[1,4,4]$.
\end{placeholder}

\begin{placeholder}[best-seed defect quotient]
\label{ph:033}
Promote the best-seed defect statement usually referred to as \textup{033} into
one manuscript-order theorem with explicit odd-$m$ quantifier:
prove that for the seed $[2,2,1]/[1,4,4]$, the defect quotient has exactly
four changed source families and four hole families, that $L3$, $R2$, and $R3$
repair directly, and that the unique unresolved channel is
$R1\to \HLone$ with
\[
\HLone=
\{(q,w,u,\Theta)=(m-1,m-1,u,2): u\neq 2\}.
\]
\end{placeholder}

\begin{theorem}[the remaining first-principles content of $M1$]
\label{thm:M1-equals-two-front-end-promotions}
For a fully first-principles manuscript-order proof, the old $M1$ placeholder
is equivalent to the conjunction of Placeholders~\ref{ph:032}
and~\ref{ph:033}.
No additional downstream channel-reduction argument remains once those two
front-end promotions are supplied.
\end{theorem}

\begin{proof}
If Placeholders~\ref{ph:032} and~\ref{ph:033} are proved, then
Theorem~\ref{thm:M1-from-front-end} applies and yields the desired channel
reduction.
Conversely, any first-principles proof of the old $M1$ statement must identify
which front-end seed is the live seed and must show that, within that seed, all
channels other than $R1\to \HLone$ are already closed. Those are exactly the
contents of Placeholders~\ref{ph:032} and~\ref{ph:033}.
\end{proof}

\section{Pilot verification already frozen by the saved defect tables}

Although the present note does not rebuild the front-end defect quotient from
scratch, the saved pilot tables still matter because they show that the exact
front-end statement being packaged here is not vague.
For the pilot moduli saved in the existing defect-template bundle, the defect
summary records exactly the following theorem-shaped content.

\begin{center}
\begin{tabular}{@{}cccc@{}}
\toprule
$m$ & direct repairs & unresolved channel & shortest unary lower bound \\
\midrule
$5$ & $L3,R2,R3$ & $R1\to \HLone$ & $50$ \\
$7$ & $L3,R2,R3$ & $R1\to \HLone$ & $196$ \\
$9$ & $L3,R2,R3$ & $R1\to \HLone$ & $486$ \\
\bottomrule
\end{tabular}
\end{center}

The value of this pilot table is not that it proves the odd-$m$ theorem by
itself. The value is that it fixes the exact front-end statement that must be
promoted: the target is not an exploratory search claim anymore, but the
specific channel theorem isolated above.

\section{Bottom line}

The old $M1$ placeholder is now much sharper.

\begin{enumerate}[label=\textup{(\arabic*)},leftmargin=2.5em]
\item Inside the current accepted front-end package, $M1$ is already closed.
\item The actual theorem proved here is the explicit reduction
      Theorem~\ref{thm:M1-from-front-end}.
\item If one still wants a completely first-principles manuscript-order proof,
      the remaining front-end externality is no longer a vague
      ``channel-classification theorem'' but exactly the pair of front-end
      promotions \textup{032} and \textup{033}.
\end{enumerate}

So after this note, the manuscript should no longer keep one large
undifferentiated $M1$ placeholder.
It should either:
\begin{enumerate}[label=\textup{(\alph*)},leftmargin=2.5em]
\item accept the current front-end packages and mark $M1$ closed; or
\item replace the old $M1$ placeholder by the two explicit front-end
      placeholders above.
\end{enumerate}

\end{document}
